\chapter{ضرب الأعداد الأكبر}\label{ch:g4:mult}

في السنة الماضية ضربنا في رقم واحد (\cref{ch:g3:mult})؛ وهذا
الفصل يضرب في أعداد من رقمين --- ويشرح «السطر المزاح» الغامض في
طريقة الأعمدة، وما هو إلا تفكيك إلى عشرات
وآحاد.

\section{الضرب في العشرات}

\begin{proposition}[الضرب في 10 و20 و300\,\dots]\label{prop:g4:mult:tens}
الضرب في $10$ يُلحق صفرًا، وفي $100$ صفرين. وللضرب
في $20$، اِضرب في $2$ ثم في $10$:
\[
36 \times 20 = 36 \times 2 \times 10 = 72 \times 10 = 720 .
\]
وكذلك $14 \times 300 = 14 \times 3 \times 100 = 4\,200$.
\end{proposition}
\admitted

\section{طريقة الأعمدة برقمين}

\begin{example}[لماذا السطر المزاح؟]\label{ex:g4:mult:why}
لحساب $23 \times 12$، فكّك $12$ إلى $10 + 2$:
\[
23 \times 12 = 23 \times 2 + 23 \times 10 = 46 + 230 = 276 .
\]
وطريقة الأعمدة تكتب هذين الجداءين بالضبط، أحدهما تحت
الآخر --- والثاني مزاح إلى اليسار لأنه عدد من
\emph{العشرات}:
\[
\begin{array}{r}
2\,3 \\
\times\ 1\,2 \\
\hline
4\,6 \\
2\,3\,\cdot \\
\hline
2\,7\,6
\end{array}
\]
(والنقطة تعلّم مرتبة الآحاد في السطر المزاح، وكثيرًا ما تُترك
فارغة).
\end{example}

\begin{omfigure}
\begin{tikzpicture}[scale=0.42]
  \fill[omDef!18] (0,0) rectangle (20,10);
  \fill[omThm!18] (0,10) rectangle (20,12);
  \fill[omDef!34] (20,0) rectangle (23,10);
  \fill[omThm!34] (20,10) rectangle (23,12);
  \draw[thick] (0,0) rectangle (23,12);
  \draw[thick] (20,0) -- (20,12);
  \draw[thick] (0,10) -- (23,10);
  \node[font=\small] at (10,5) {$20 \times 10 = 200$};
  \node[font=\small] at (10,11) {$20 \times 2 = 40$};
  \node[font=\small, rotate=90] at (21.5,5) {$3 \times 10 = 30$};
  \node[font=\small] at (21.5,11) {$6$};
  \node[below, font=\small] at (10,0) {$20$};
  \node[below, font=\small] at (21.5,0) {$3$};
  \node[left, font=\small] at (0,5) {$10$};
  \node[left, font=\small] at (0,11) {$2$};
\end{tikzpicture}

{\small صورة المستطيل \omterm{def:g2:mult:def}{للجداء} $23 \times 12$: أربعة \omterm{def:g2:mult:def}{جداءات} سهلة،
$200 + 40 + 30 + 6 = 276$. وطريقة الأعمدة تجمعها في سطرين:
$46$ (سطر «$\times 2$») و $230$ (سطر «$\times 10$»).}
\end{omfigure}

\begin{method}[الضرب بالأعمدة في عدد من رقمين]\label{met:g4:mult:column}
لحساب $457 \times 36$:
\begin{enumerate}
  \item السطر الأول: $457 \times 6 = 2\,742$؛
  \item والسطر الثاني: $457 \times 3$ \emph{عشرات} $= 1\,371$
        عشرة --- اُكتب $1\,371$ مزاحًا مرتبة واحدة إلى اليسار؛
  \item ثم اِجمع السطرين: $2\,742 + 13\,710 = 16\,452$؛
  \item وقدِّر للتحقّق: $457 \times 36 \approx 500 \times 36
        = 18\,000$؟ والأفضل: $450 \times 36$ نحو $16\,000$ ---
        فالنتيجة معقولة.
\end{enumerate}
\[
\begin{array}{r}
4\,5\,7 \\
\times\ \ \,3\,6 \\
\hline
2\,7\,4\,2 \\
1\,3\,7\,1\,\cdot \\
\hline
1\,6\,4\,5\,2
\end{array}
\]
\end{method}

\begin{example}[حيل ذهنية]\label{ex:g4:mult:tricks}
\begin{itemize}
  \item \emph{فكّك بالطريقة الودودة}: $18 \times 11 = 18 \times 10
        + 18 = 198$.
  \item \emph{اِستعمل قريبًا من المئة}: $25 \times 99 = 25 \times 100 -
        25 = 2\,475$.
  \item \emph{ضاعِف ونصِّف}: $16 \times 35 = 8 \times 70 =
        560$ (فتنصيف أحد العاملين وتضعيف الآخر يحفظ
        \omterm{def:g2:mult:def}{الجداء}).
\end{itemize}
\end{example}

\section{الضرب في المسائل}

\begin{example}\label{ex:g4:mult:problem}
تطلب مدرسة $28$ علبة فيها $145$ ورقة لكل علبة.
\begin{enumerate}
  \item العملية: $145 \times 28$.
  \item التقدير: $150 \times 30 = 4\,500$، وهو أكثر قليلًا في
        العاملين معًا.
  \item الأعمدة: $145 \times 8 = 1\,160$؛ و $145 \times 2$ عشرة
        $= 2\,900$؛ \omterm{def:g1:addition:def}{والمجموع} $4\,060$.
  \item والجملة: تستلم المدرسة $4\,060$ ورقة.
\end{enumerate}
\end{example}

\section{تمارين}

\begin{exercise}[$\star$]\label{exo:g4:mult:1}
احسب ذهنيًّا: $47 \times 10$؛ \quad $23 \times 100$؛ \quad
$36 \times 20$؛ \quad $15 \times 300$؛ \quad $42 \times 50$.
\end{exercise}

\begin{exercise}[$\star$]\label{exo:g4:mult:2}
احسب $34 \times 21$ بالتفكيك كما في \cref{ex:g4:mult:why}
($34 \times 21 = 34 \times 20 + 34$)، ثم تحقّق بالأعمدة.
\end{exercise}

\begin{exercise}[$\star$]\label{exo:g4:mult:3}
احسب بالأعمدة: $63 \times 24$؛ \quad $85 \times 47$؛ \quad
$126 \times 53$.
\end{exercise}

\begin{exercise}[$\star$]\label{exo:g4:mult:4}
احسب بالأعمدة، مع تقدير أولًا: $208 \times 34$؛ \quad
$517 \times 62$.
\end{exercise}

\begin{exercise}[$\star$]\label{exo:g4:mult:5}
احسب بحيلة من \cref{ex:g4:mult:tricks}:
$45 \times 11$؛ \quad $32 \times 99$؛ \quad $24 \times 45$
(ضاعِف ونصِّف --- مرتين إن شئت).
\end{exercise}

\begin{exercise}[$\star$]\label{exo:g4:mult:6}
في سينما $26$ صفًّا من $18$ مقعدًا. فكم مقعدًا فيها كلها؟
\end{exercise}

\begin{exercise}[$\star$]\label{exo:g4:mult:7}
تحمل شاحنة $32$ منصّة فيها $48$ صندوقًا لكل منصّة. فكم صندوقًا؟
أعطِ تقديرًا أولًا، ثم الجواب الدقيق.
\end{exercise}

\begin{exercise}[$\star$]\label{exo:g4:mult:8}
جِد الخطأ: تحسب سلمى $57 \times 23$ هكذا:
$57 \times 3 = 171$ و $57 \times 2 = 114$، ثم تجمع
$171 + 114 = 285$. والجواب الصحيح هو $1\,311$. فماذا
نسيت؟
\end{exercise}

\begin{exercise}[$\star\star$]\label{exo:g4:mult:9}
يغرس بستانيّ $14$ صفًّا من $35$ زهرة تُوليب و $12$ صفًّا من $28$
نرجسة. فكم زهرة في المجموع؟ (جداءان، ثم \omterm{def:g1:addition:def}{مجموع}.)
\end{exercise}

\begin{exercise}[$\star\star$]\label{exo:g4:mult:10}
في كل يوم، يستعمل مخبز $67$ كغ من الدقيق. فكم يقارب ذلك من
الدقيق في شهر من $30$ يومًا --- قدِّر بأعداد مقرَّبة، ثم احسب
بالضبط.
\end{exercise}

\begin{exercise}[$\star\star$]\label{exo:g4:mult:11}
اِشرح، بصورة مستطيل كصورة هذا الفصل، لماذا يمكن حساب
$15 \times 12$ على صورة $15 \times 10 + 15 \times 2$.
ثم اِبتكر تفكيكًا آخر \omterm{def:g2:mult:def}{للجداء} $15 \times 12$ ينجح أيضًا.
\end{exercise}
